% Generated draft: requires mathematical and source review.
\documentclass[11pt]{article}
\usepackage[margin=25mm]{geometry}
\usepackage{fontspec}
\setmainfont{Latin Modern Roman}
\usepackage{amsmath,amssymb,mathtools,unicode-math}
\setmathfont{Latin Modern Math}
\usepackage{enumitem,needspace,xurl,hyperref}
\setlength{\parindent}{0pt}
\setlength{\parskip}{5pt}
\setlength{\emergencystretch}{3em}
\newcommand{\N}{\mathbb{N}}
\newcommand{\Z}{\mathbb{Z}}
\newcommand{\Q}{\mathbb{Q}}
\newcommand{\R}{\mathbb{R}}
\newcommand{\C}{\mathbb{C}}
\newcommand{\F}{\mathbb{F}}
\newcommand{\A}{\mathbb{A}}
\newcommand{\PP}{\mathbb{P}}
\newcommand{\E}{\mathbb{E}}
\newcommand{\eps}{\varepsilon}
\newcommand{\dd}{\,\mathrm{d}}
\begin{document}
\section*{UnsolvedMath ID: 30002427}
\textbf{Source problem code:} catalogue entry
\subsection{Definitions and mathematical statement}
Let \(N\) be a positive integer and set \(R_N=\Z/N\Z\). A subgroup \(G\leq\operatorname{GL}_2(R_N)\) has \emph{full determinant} if
\[
\det(G)=R_N^{\times}.
\]
For the purposes of this draft, use the usual compactified coarse modular curve convention: \(X_G/\Q\) is the smooth projective modular curve obtained by compactifying the coarse moduli curve associated with \(G\). Over \(\overline{\Q}\), its noncuspidal moduli points represent elliptic curves equipped with a \(G\)-orbit of ordered bases of their \(N\)-torsion, up to isomorphism. The model over \(\Q\) is the standard modular model associated with \(G\). At \(N=1\), use the usual level-one curve. Write
\[
X_G(\Q)=\operatorname{Hom}_{\Q}(\operatorname{Spec}\Q,X_G).
\]
This denotes rational points of the coarse curve; no identification with individual level structures defined over \(\Q\) is asserted. The source does not specify compactification or the treatment of \(-I\). Under the convention used here, cusps are included, and replacing \(G\) by \(\langle G,-I\rangle\) leaves the coarse modular curve unchanged. These conventions require confirmation against the intended source formulation.

For a prime \(p\), identify \(R_p\) with \(\F_p\). A standard Borel subgroup and a standard split Cartan subgroup are, respectively,
\[
B_p=\left\{\begin{pmatrix}a&b\\0&d\end{pmatrix}:a,d\in\F_p^{\times},\ b\in\F_p\right\},
\qquad
C_{\mathrm{s}}(p)=\left\{\begin{pmatrix}a&0\\0&d\end{pmatrix}:a,d\in\F_p^{\times}\right\}.
\]
Choose an \(\F_p\)-basis of \(\F_{p^2}\). Multiplication by nonzero elements of \(\F_{p^2}\), expressed in this basis, gives a nonsplit Cartan subgroup
\[
C_{\mathrm{ns}}(p)=\{m_\alpha:\alpha\in\F_{p^2}^{\times}\}\leq\operatorname{GL}_2(\F_p),
\qquad m_\alpha(x)=\alpha x.
\]
The basis choice changes this subgroup only by conjugacy. For any subgroup \(C\leq\operatorname{GL}_2(\F_p)\), define its normalizer by
\[
\operatorname{Norm}(C)=\{g\in\operatorname{GL}_2(\F_p):gCg^{-1}=C\}.
\]
Use the notation \(X_{\mathrm{s}}^+(p)=X_{\operatorname{Norm}(C_{\mathrm{s}}(p))}\) and \(X_{\mathrm{ns}}^+(p)=X_{\operatorname{Norm}(C_{\mathrm{ns}}(p))}\).

\textbf{Problem.} For every positive integer \(N\) and every full-determinant subgroup \(G\leq\operatorname{GL}_2(\Z/N\Z)\), describe the complete set \(X_G(\Q)\). In particular, carry out this determination for every prime \(p\) when \(G\) is a Borel subgroup, the normalizer of a split Cartan subgroup, or the normalizer of a nonsplit Cartan subgroup of \(\operatorname{GL}_2(\F_p)\), including their conjugates. The source states no restriction on the prime. It does not prescribe the form of a description, such as a parametrization, an explicit list, or a uniform classification theorem; that intended output remains unresolved.
\subsection{Short English statement}
Determine all rational points on modular curves associated with full-determinant matrix subgroups, with particular emphasis on prime-level Borel and Cartan-normalizer curves. This draft explicitly uses compactified coarse modular curves, so rational cusps are included; the source does not specify these conventions. The supplied label is partially solved, and the accompanying literature assessment requires independent review.
\subsection{Sources}
\begin{itemize}
\item Local UnsolvedMath registry, permanent ID \(30002427\), source problem code \texttt{catalogue entry}, title \emph{Rational Points on Cartan Modular Curves}, in the Oberwolfach Reports Open Problems collection. The supplied record gives proposed year 2013 and names no individual proposer.
\item Supplied report attribution: \emph{Analytic Number Theory (2014)}, Oberwolfach Reports, DOI \texttt{10.4171/owr/2013/51}, pp.\ 2963--3037. \url{https://doi.org/10.4171/owr/2013/51}. The record describes Oberwolfach Reports as documenting mathematical workshops and their open-problem discussions. No precise problem page or named contributor is supplied.
\item Supplied literature reference, labelled ``Bilu--Parent split-Cartan theorem'': \url{https://doi.org/10.4007/annals.2011.173.1.13}.
\item Supplied literature reference, labelled ``Bilu--Parent--Rebolledo rational points'': \url{https://doi.org/10.5802/aif.2781}.
\item Supplied literature reference, labelled ``quadratic Chabauty computations'': \url{https://arxiv.org/abs/2101.01862}.
\item Literature assessment dated exactly \texttt{2026-08-21}: the supplied curation labels the problem partially solved and states that rational points on split-Cartan and related Atkin--Lehner curves are classified in broad ranges and exceptional levels, including level 13, while the general rational-point problem for \(X_{\mathrm{ns}}^+(p)\) remains open. This is supplied dated curation derived from the OWR material, not an independently verified literature assessment; the record itself requests independent verification of time-sensitive status and cited claims.
\item The supplied statement-status label is \texttt{corrected\_verified}. Its accompanying note says that notation, typography, and wording were normalized without changing the mathematical content. This records the supplied editorial assessment and does not establish independent mathematical verification of this draft.
\end{itemize}
\end{document}
